% Formal Learning Theory Kernel: Blueprint v1
% Content source. This file is included by print.tex and web.tex;
% it should not carry \documentclass or \begin{document}.
%
% Scope: core story of the kernel. Five-way PAC equivalence with the
% NullMeasurableSet refinement, Littlestone characterization, Gold
% theorem and mind-change characterization, three-paradigm separation,
% Borel-analytic measurability refinement. Compression, MBL monad, and
% PAC-Bayes are deferred to blueprint v2.

\chapter{Introduction}\label{chap:intro}

This document is the mathematical blueprint for a Lean 4 formalization
of the fundamental theorem of statistical learning. It tracks the core
story of the kernel: the five-way equivalence for PAC learning with a
measurability refinement that weakens the standard Borel hypothesis,
the Littlestone characterization for online learning, Gold's theorem
and the mind-change characterization for Gold-style learning, the
three-paradigm separation, and the Borel-analytic separation theorem
that witnesses the strict weakening of the measurability hypothesis.

The formalization targets Lean 4 \texttt{v4.29.0-rc6} against mathlib4
pinned at commit \texttt{fde0cc5}. The repository sits at
\url{https://github.com/Zetetic-Dhruv/formal-learning-theory-kernel},
the per-module API reference at
\url{https://zetetic-dhruv.github.io/formal-learning-theory-kernel/}.
The kernel contains 354 machine-checked theorems in 21{,}522 lines with
zero sorry. Every statement in this blueprint carrying a
\texttt{\textbackslash lean\{\}} annotation corresponds to a proved
declaration in the kernel.

The audience is learning theorists who do not read Lean 4. Informal
statements are self-contained mathematical prose. The
\texttt{\textbackslash lean\{\}} hyperlinks are there for formalization
readers who want to drill into the precise Lean counterparts of each
statement. Both audiences are served by the same document.

Compression via approximate minimax, the measurable batch learner
monad, and PAC-Bayes infrastructure are deferred to blueprint v2,
together with the Baxter multi-task base case and the extended
criteria family (robust PAC, PAC-Bayes, RKHS statistical layer).
Blueprint v2 will be published as a separate document once v1 is
complete.

\chapter{Primitives}\label{chap:primitives}

\section{Concept and concept class}\label{sec:concept-class}

\begin{definition}[Concept and concept class]\label{def:concept-class}
  \lean{Concept, ConceptClass}
  \leanok
  Fix a domain $X$ and a label set $Y$. A \emph{concept} is a function
  $h : X \to Y$. A \emph{concept class} is a set $C \subseteq Y^X$ of
  concepts.
\end{definition}

Throughout this blueprint we work mostly with the Boolean label set
$Y = \{0, 1\}$, identifying $\{0,1\}^X$ with the power set of $X$. All
definitions extend to arbitrary $Y$ unless stated otherwise, and the
kernel formalizes the general case.

\section{Batch learner}\label{sec:batch-learner}

\begin{definition}[Batch learner]\label{def:batch-learner}
  \lean{BatchLearner}
  \leanok
  A \emph{batch learner} over $(X, Y)$ is a bundled structure carrying a
  hypothesis set $H \subseteq Y^X$, a learning function
  $\mathcal{A}$ that takes a finite training sample
  $S : \mathrm{Fin}\,m \to X \times Y$ and returns a hypothesis in $H$,
  and a proof that the output lies in $H$.
\end{definition}

\begin{definition}[Realizable setting]\label{def:realizable}
  The learning problem for a concept class $C$ is \emph{realizable} if
  the target concept $c$ lies in $C$ and the sample is drawn i.i.d.\
  from a distribution $D$ with labels $c(x_i)$. Unless otherwise noted,
  PAC learnability statements in this blueprint are in the realizable
  setting. The kernel also formalizes the agnostic variant and the
  extended criteria in the extended theorems module.
\end{definition}

\section{Sample, loss, and error}\label{sec:sample-loss}

\begin{definition}[I.i.d.\ sample]\label{def:iid-sample}
  For a distribution $D$ on $X$ and a target concept $c$, an
  \emph{i.i.d.\ sample} of size $m$ is a tuple
  $S = ((x_1, c(x_1)), \dots, (x_m, c(x_m)))$ with each $x_i \sim D$
  drawn independently.
\end{definition}

\begin{definition}[Zero-one loss]\label{def:zero-one-loss}
  For $h, c : X \to \{0,1\}$ and $x \in X$, the \emph{zero-one loss} is
  $\ell(h, c, x) = \mathbf{1}[h(x) \neq c(x)]$.
\end{definition}

\begin{definition}[Empirical and true error]\label{def:errors}
  For a hypothesis $h$, a target $c$, a distribution $D$, and a sample
  $S$ of size $m$:
  \begin{align*}
    \err_S(h) &= \frac{1}{m} \sum_{i=1}^{m} \mathbf{1}[h(x_i) \neq c(x_i)]
      && \text{(empirical error)}, \\
    \err_D(h) &= \mathbb{P}_{x \sim D}\!\left[ h(x) \neq c(x) \right]
      && \text{(true error)}.
  \end{align*}
\end{definition}

\chapter{PAC Characterization}\label{chap:pac}

\section{VC dimension, shattering, and the growth function}\label{sec:vc-dim}

\begin{definition}[Shattering]\label{def:shatters}
  \lean{Shatters}
  \leanok
  A finite set $S \subseteq X$ is \emph{shattered} by a concept class
  $C \subseteq \{0,1\}^X$ if for every function $f : S \to \{0,1\}$
  there exists $h \in C$ with $h \restriction_S = f$. Equivalently, $C$
  realizes every Boolean labeling of $S$.
\end{definition}

\begin{definition}[VC dimension]\label{def:vcdim}
  The \emph{VC dimension} of $C$, written $\VCdim(C)$, is the supremum
  of cardinalities of finite sets shattered by $C$:
  \[
    \VCdim(C) = \sup\{ |S| : S \subseteq X \text{ finite and shattered by } C \}.
  \]
  If the supremum is not finite, we set $\VCdim(C) = \infty$.
\end{definition}

\begin{definition}[Growth function]\label{def:growth}
  The \emph{growth function} $\Pi_C(m) : \N \to \N$ is
  \[
    \Pi_C(m) = \sup_{\substack{S \subseteq X \\ |S| = m}}
      \bigl|\, \{ h \restriction_S : h \in C \}\, \bigr|.
  \]
  It counts the maximum number of distinct Boolean labelings of an
  $m$-element subset that the class can realize.
\end{definition}

\section{Sauer-Shelah lemma}\label{sec:sauer-shelah}

The Sauer-Shelah lemma is the combinatorial heart of the uniform
convergence argument. Without it, the growth function on a class with
finite VC dimension could grow exponentially in the sample size; with
it, the growth is polynomial of degree at most $\VCdim(C)$, and the
union bound over growth-function-many effective hypotheses closes the
PAC proof.

\begin{lemma}[Sauer-Shelah]\label{lem:sauer-shelah}
  \lean{sauer_shelah}
  \leanok
  \uses{def:vcdim, def:growth}
  Let $C$ be a concept class with $\VCdim(C) = d < \infty$. Then for
  every $m \geq 0$,
  \[
    \Pi_C(m) \leq \sum_{i=0}^{d} \binom{m}{i}.
  \]
  In particular, $\Pi_C(m) = O(m^d)$ as $m \to \infty$.
\end{lemma}

\begin{proof}[Proof sketch]
  The classical proof is a downward induction on VC dimension via a
  shift operator on the restriction $C \restriction_S$. For each element
  $x \in S$, partition the restrictions into those distinguished at $x$
  and those that are not; the distinguished subclass has VC dimension
  at most $d - 1$, and a short double-counting argument with Pascal's
  identity closes the induction. The kernel formalizes the quantitative
  variant with explicit binomial-sum bounds as
  \lean{sauer_shelah_quantitative}, and the exponential upper bound
  $\Pi_C(m) \leq (em/d)^d$ used downstream by the symmetrization
  argument as \lean{sauer_shelah_exp_bound}.
\end{proof}

\section{The five-way equivalence}\label{sec:fundamental-theorem}

This section is the centerpiece of the blueprint. The fundamental
theorem bundles five a priori distinct characterizations of PAC
learnability into a single equivalence class, and this kernel adds a
measurability refinement to the statement: the Borel hypothesis that
appears in the standard treatment can be weakened to a
null-measurability hypothesis on the uniform convergence bad event.
The weakening is strict, witnessed by the class constructed in
Chapter~\ref{chap:measurability}.

\begin{definition}[Uniform convergence]\label{def:uniform-convergence}
  \lean{HasUniformConvergence}
  \leanok
  A concept class $C$ has the \emph{uniform convergence property} over
  all distributions on $X$ if for every $\eps, \delta \in (0, 1)$ there
  exists a sample complexity $m_0(\eps, \delta)$, independent of the
  distribution, such that for every distribution $D$ on $X$,
  \[
    \mathbb{P}_{S \sim D^{m_0}}
      \!\left[ \sup_{h \in C} \bigl| \err_S(h) - \err_D(h) \bigr| > \eps \right]
      < \delta.
  \]
\end{definition}

\begin{definition}[PAC learnability]\label{def:pac-learnable}
  \lean{PACLearnable}
  \leanok
  A concept class $C$ is \emph{PAC learnable} if there exist a batch
  learner $\mathcal{A}$ and a sample complexity
  $m_0 : (0, 1)^2 \to \N$, both independent of the target distribution,
  such that for every $\eps, \delta \in (0, 1)$, every distribution $D$,
  and every realizable target $c \in C$,
  \[
    \mathbb{P}_{S \sim D^{m_0(\eps, \delta)}}
      \!\left[ \err_D(\mathcal{A}(S)) > \eps \right] < \delta.
  \]
\end{definition}

\begin{theorem}[Five-way equivalence for PAC learning]\label{thm:fundamental-theorem}
  \lean{fundamental_theorem}
  \leanok
  \uses{def:vcdim, def:uniform-convergence, def:pac-learnable,
        lem:sauer-shelah, def:well-behaved-vc-meas-target}
  Let $X$ be an infinite domain and let $C \subseteq \{0,1\}^X$ be a
  concept class carrying \texttt{WellBehavedVCMeasTarget}
  (Definition~\ref{def:well-behaved-vc-meas-target}). The following are
  equivalent.
  \begin{enumerate}[label=(\roman*)]
    \item $\VCdim(C) < \infty$.
    \item $C$ has the uniform convergence property over all
      distributions on $X$.
    \item Empirical risk minimization over $C$ is a PAC learner.
    \item $C$ is PAC learnable.
    \item Every realizable distribution admits a sample complexity
      bound polynomial in $1/\eps$ and $1/\delta$.
  \end{enumerate}
\end{theorem}

\begin{proof}[Proof sketch]
  Forward direction (i) $\Rightarrow$ (ii). Sauer-Shelah
  (Lemma~\ref{lem:sauer-shelah}) gives polynomial growth of the growth
  function $\Pi_C(2m)$. Symmetrization introduces a ghost sample and
  rewrites the supremum over $C$ as a supremum over the finitely many
  effective labelings of the combined sample (see
  Section~\ref{sec:symmetrization-nullmeas}). Hoeffding's inequality
  bounds the deviation at each fixed effective labeling, and the union
  bound over $\Pi_C(2m)$ such labelings produces the uniform
  convergence rate. The symmetrization step uses the
  \texttt{WellBehavedVCMeasTarget} hypothesis to ensure the ghost-gap
  bad event is null-measurable, which is the exact regularity the
  argument requires.

  Chain (ii) $\Rightarrow$ (iii) $\Rightarrow$ (iv) $\Rightarrow$ (v).
  Uniform convergence implies ERM soundness in the realizable setting:
  the empirical risk minimizer has zero empirical error, and the
  uniform convergence bound transfers zero empirical error to
  vanishing true error. The sample complexity bound from (ii) carries
  over to PAC learnability. The polynomial rate (v) is extracted from
  the explicit $m_0(\eps, \delta)$ produced by the union bound: it is
  polynomial in $1/\eps$ and $\log(1/\delta)$, and therefore polynomial
  in $1/\eps$ and $1/\delta$ as claimed.

  Backward direction (v) or (iv) $\Rightarrow$ (i). Contrapositive: if
  $\VCdim(C) = \infty$, then for every $m$ there is a shattered set
  $S_m$ of size $|S_m| > m$. A no-free-lunch argument on $S_m$
  constructs, for any learner and any sample-complexity bound $m$, a
  realizable target and a distribution on $S_m$ under which the
  learner's expected error exceeds a fixed constant with constant
  probability. Taking $|S_m|$ large enough forces the error above any
  prescribed $\eps$ with probability at least $\delta$, contradicting
  PAC learnability.
\end{proof}

\subsection{The NullMeasurableSet refinement}\label{sec:symmetrization-nullmeas}

The proof of (i) $\Rightarrow$ (ii) above goes through a
symmetrization argument that rewrites the bad event
\[
  \mathrm{Bad}_1(C, m, \eps)
    = \left\{ S \in X^m :
      \sup_{h \in C} \bigl| \err_S(h) - \err_D(h) \bigr| > \eps \right\}
\]
in terms of a ghost sample $S'$ and the two-sided comparison event
\[
  \mathrm{Bad}_2(C, m, \eps)
    = \left\{ (S, S') \in X^m \times X^m :
      \sup_{h \in C} \bigl| \err_{S'}(h) - \err_S(h) \bigr| > \eps \right\}.
\]
The probability bound on $\mathrm{Bad}_2$ requires it to be measurable
with respect to the product measure; the symmetrization step is where
regularity of the class enters the proof.

\begin{definition}[WellBehavedVCMeasTarget]\label{def:well-behaved-vc-meas-target}
  \lean{WellBehavedVCMeasTarget}
  \leanok
  A concept class $C$ over a measurable space $(X, \Sigma)$ satisfies
  \emph{WellBehavedVCMeasTarget} if the two-sided ghost-gap bad event
  $\mathrm{Bad}_2(C, m, \eps)$ is a \emph{null-measurable} subset of
  $(X^m \times X^m,\, \Sigma^{\otimes 2m})$ with respect to every
  product probability measure $D^{\otimes 2m}$ arising from a
  distribution $D$ on $X$. Null-measurability means that
  $\mathrm{Bad}_2$ sits in the completion of the product
  $\sigma$-algebra.
\end{definition}

\begin{definition}[KrappWirth well-behaved]\label{def:krapp-wirth-well-behaved}
  \lean{KrappWirthWellBehaved}
  \leanok
  A concept class $C$ satisfies the \emph{KrappWirth well-behaved}
  hypothesis of Krapp and Wirth (2024) if the two-sided ghost-gap bad
  event $\mathrm{Bad}_2(C, m, \eps)$ is a Borel subset of
  $X^m \times X^m$ in the product Borel structure.
\end{definition}

\begin{proposition}[KrappWirth implies WellBehavedVCMeasTarget]\label{prop:krapp-wirth-implies-wb-vc}
  \lean{KrappWirthWellBehaved.toWellBehavedVC}
  \leanok
  \uses{def:well-behaved-vc-meas-target, def:krapp-wirth-well-behaved}
  Every Borel set is null-measurable with respect to any completion of
  a Borel probability measure. Consequently, any concept class carrying
  \texttt{KrappWirthWellBehaved} also carries
  \texttt{WellBehavedVCMeasTarget}.
\end{proposition}

\begin{proof}
  Borel sets lie in the product $\sigma$-algebra by definition, hence
  in its completion with respect to any measure. The inclusion
  $\mathcal{B} \subseteq \mathcal{B}_\mu^\ast$ (Borel into
  completion) is immediate.
\end{proof}

\begin{remark}[Why the refinement matters]\label{rem:why-refinement}
  The standard treatment in Krapp and Wirth (2024) states the
  fundamental theorem under the Borel hypothesis on the ghost-gap bad
  event, matching Definition~\ref{def:krapp-wirth-well-behaved}. This
  kernel observes that the symmetrization step uses only
  null-measurability of the bad event, not Borelness, and that the
  resulting theorem (Theorem~\ref{thm:fundamental-theorem}) is strictly
  stronger: it applies to concept classes whose ghost-gap bad event is
  null-measurable without being Borel.
  Chapter~\ref{chap:measurability} exhibits a concept class whose bad
  event sits precisely in the gap between Borel and null-measurable,
  certifying that the refinement is not cosmetic.
\end{remark}

\section{Forward pointer to the measurability layer}\label{sec:pac-pointer-forward}

The fundamental theorem (Theorem~\ref{thm:fundamental-theorem}) uses
the \texttt{WellBehavedVCMeasTarget} hypothesis as a black box at this
stage. Chapter~\ref{chap:measurability} opens that black box: it
constructs the Borel-parameterized setting in which the ghost-gap bad
event is always \emph{analytic} (a strictly larger class than Borel),
proves that analytic sets are null-measurable with respect to every
Borel probability measure via Choquet capacitability, and exhibits a
concept class (a singleton class over an analytic non-Borel set) whose
bad event is null-measurable but not Borel. This witnesses the strict
weakening announced in Remark~\ref{rem:why-refinement}.

% Chapters 3 through 7 deferred to later implementation sessions.
% Planned: Chapter 3 Online characterization, Chapter 4 Gold
% characterization, Chapter 5 Three-paradigm separations, Chapter 6
% Measurability layer, Chapter 7 Closing notes. See blueprint/PLAN.md
% for scope and structure.
